\specialchapter{Prologue}{Two parameters}{{2/2}{2}{\tfrac12}{a^2b^2}}
\sltfloatprefix{P}
\begin{refsection}[chapters/00-prologue.bib]

Take two numbers $a$ and $b$, multiply them, and use the product as the slope of a line through the origin:
\[
y=abx+\varepsilon,\qquad x\sim\Normal(0,1),\quad\varepsilon\sim\Normal(0,1).
\]
This is the smallest model in this book, and the one we return to in every chapter. Fit it, in the Bayesian way, to data that are pure noise: $n$ pairs $(x_i,y_i)$ drawn with $y_i=\varepsilon_i$, so that the truth is the slope $0$. Put a uniform prior on the square $[-1,1]^2$, and look at the posterior distribution of $(a,b)$.

A statistician trained on regular models expects to see a small Gaussian blob. The model has two parameters; there is a maximum-likelihood estimate $(\hat a,\hat b)$; the posterior should concentrate around it at the rate $n^{-1/2}$, with a covariance given by the inverse Fisher information; and the logarithm of the marginal likelihood should be the maximized log-likelihood minus $\tfrac d2\log n=\log n$, up to a bounded term. That is what the Bayesian information criterion of \textcite{Schwarz1978} says, and it is what every textbook proof of the Laplace approximation delivers.

None of it happens. \Cref{fig:prologue:draws} shows draws from the posterior for three sample sizes. They do not form a blob. They crowd onto the cross $ab=0$, the set of all parameters that give the true slope, and they spread along both arms. As $n$ grows, the arms get thinner but no shorter, and the draws do not converge to a point. There is no maximum-likelihood point to converge to, either: the likelihood depends on $(a,b)$ only through the product $c=ab$, and every point of the hyperbola $ab=\hat c$ maximizes it.

\begin{figure}[t]
\centering
\begin{tikzpicture}
\begin{groupplot}[group style={group size=3 by 1,horizontal sep=4mm,yticklabels at=edge left},
  splot,width=0.39\linewidth,height=0.39\linewidth,xmin=-1,xmax=1,ymin=-1,ymax=1,
  xtick={-1,0,1},ytick={-1,0,1},xlabel={$a$},title style={yshift=-1ex}]
\nextgroupplot[ylabel={$b$},title={$n=10^2$}]
\addplot[draws] table[x=a,y=b] {figures/data/00-prologue-draws-n2.dat};
\nextgroupplot[title={$n=10^4$}]
\addplot[draws] table[x=a,y=b] {figures/data/00-prologue-draws-n4.dat};
\nextgroupplot[title={$n=10^6$}]
\addplot[draws] table[x=a,y=b] {figures/data/00-prologue-draws-n6.dat};
\end{groupplot}
\end{tikzpicture}
\caption{Draws from the posterior of $y=abx$ when the truth is the slope $0$, with the uniform prior on $[-1,1]^2$: $1500$ exact draws for one data set at each sample size (data seeds $102$, $104$, $106$). The draws first fix the product $c=ab$ from its one-dimensional posterior, then $a$ given $c$; \cref{ch:singular} derives the method. The posterior collapses onto the cross $ab=0$ but not onto a point.}\label{fig:prologue:draws}
\end{figure}

The marginal likelihood misbehaves too. Write $F_n$ for its negative logarithm, the \term{free energy}\index{free energy}, and $nS_n$ for the negative log-likelihood of the truth, so that $F_n-nS_n$ is what the data charge the model for not knowing the truth in advance. \Cref{fig:prologue:free} plots the average of $F_n-nS_n$ over two hundred data sets against $n$, for our model and for the regular model $y=ax_1+bx_2$, which also has two parameters, also has the uniform prior on $[-1,1]^2$, and also has the truth $0$. The regular model pays exactly what Schwarz predicts: one unit of $\log n$ for its two parameters, plus a constant. Our model pays about half as much, and a little less again: its free energy grows like
\begin{equation}\label{eq:prologue:free}
F_n-nS_n=\tfrac12\log n-\log\log n+O_p(1).
\end{equation}
The measured slope against $\log n$ over the range of the figure is $0.39$, below $\tfrac12$ because of the term $-\log\log n$, and the bronze curve, computed from the formula that \cref{ch:coefficient} proves, runs through every measured point.

\begin{figure}[t]
\centering
\begin{tikzpicture}
\begin{axis}[splot,xmode=log,xlabel={sample size $n$},ylabel={$\E[F_n-nS_n]$},ymin=0,ymax=17.5,
  ytick={0,4,8,12},legend pos=north west]
\addplot[regularpred] table[x=n,y=regpred] {figures/data/00-prologue-free.dat};
\addplot[measured3,ebars] table[x=n,y=reg,y error plus=reghi,y error minus=reglo] {figures/data/00-prologue-free.dat};
\addplot[theory] table[x=n,y=prodpred] {figures/data/00-prologue-free.dat};
\addplot[measured,ebars] table[x=n,y=prod,y error plus=prodhi,y error minus=prodlo] {figures/data/00-prologue-free.dat};
\legend{regular prediction: $\log n+$const.,$y=ax_1+bx_2$ (measured),singular prediction: $\tfrac12\log n-\log\log n+\dots$,$y=abx$ (measured)}
\end{axis}
\end{tikzpicture}
\caption{The normalized free energy $F_n-nS_n$ of two models with two parameters each, averaged over $200$ data sets at each of nine sample sizes from $10^2$ to $10^6$, with Student-$t$ $95\%$ intervals; the truth is $0$ and the prior uniform on $[-1,1]^2$ for both. Squares: the regular model $y=ax_1+bx_2$ (Gaussian integral). Dots: the model $y=abx$ (exact quadrature in $c=ab$). Dashed: Laplace's prediction $\log n-1-\log(2\pi)+\log4$. Solid: the prediction of \cref{thm:coefficient:free} with the slowly varying factor kept, which tends to $\tfrac12\log n-\log\log n-\tfrac12+\log(4/\sqrt{2\pi})$. Script: \texttt{figures/scripts/00-prologue-free.py}.}\label{fig:prologue:free}
\end{figure}

The two numbers in \eqref{eq:prologue:free}, the coefficient $\tfrac12$ of $\log n$ and the coefficient $1$ of $\log\log n$, are the subject of this book. The first is the \term{learning coefficient}\index{learning coefficient} $\lambda$ of the model at its truth; the second is $m-1$, where $m$ is the learning coefficient's \term{multiplicity}\index{multiplicity}. For a regular model with $d$ parameters, $\lambda=\tfrac d2$ and $m=1$, and the classical formulas are the special case. For $y=abx$ at the truth $0$, $\lambda=\tfrac12$ and $m=2$: two parameters that learn like one, with a bonus.

\section*{Why two parameters learn like one}
\addcontentsline{toc}{section}{Why two parameters learn like one}

The explanation is geometric, and it fits in a paragraph. The data can only see the product $ab$, so the parameters $(a,b)$ and $(ta,b/t)$ give the same distribution for every $t\neq0$: the model folds its parameter space onto the interval of slopes. Learning the truth means locating the slope $ab$ to within about $n^{-1/2}$, and the posterior puts its mass on the parameters that do so, which are those of the thin region $\{\abs{ab}\lesssim n^{-1/2}\}$ around the cross. What matters for the free energy is the prior mass of that region. A regular model with two parameters must locate both coordinates, and the region it keeps is a disc of radius $n^{-1/2}$, of area about $n^{-1}$; the free energy pays the logarithm of the inverse, $\log n$. Our model keeps a cross-shaped region whose area is about $n^{-1/2}\log n$: the arms are $n^{-1/2}$ thick far from the centre and fatter near it. The free energy pays $\tfrac12\log n-\log\log n$. In general, the volume $V(\eps)$ of the parameters whose loss is within $\eps$ of the truth's shrinks like $\eps^\lambda(\log\frac1\eps)^{m-1}$ as $\eps\to0$, and the free energy pays $\lambda\log n-(m-1)\log\log n$. \Cref{ch:coefficient,ch:volume} make this precise.

The same geometry controls everything else a statistician asks of a model. The expected Kullback--Leibler divergence from the truth to the Bayes predictive distribution, the generalization error, is $\lambda/n$ to first order, not $\tfrac d{2n}$ (\cref{ch:generalization}). The penalties of the information criteria change accordingly (\cref{ch:waic,ch:temperature}). And the posterior does not settle near one point: when several regions of the parameter space nearly fit the data, it moves between them as $n$ grows, trading accuracy against learning coefficient (\cref{ch:phases}).

\section*{Why this matters for neural networks}
\addcontentsline{toc}{section}{Why this matters for neural networks}

A model with a cross for its set of true parameters can look like a curiosity designed to break a theorem. It is not; it is the smallest piece of a neural network. Take one hidden unit with input weight $b$ and output weight $a$. With a $\tanh$ activation its contribution is $a\tanh(bx)$, which is $abx$ to first order near $b=0$; with a ReLU it is $a\relu(bx)$, which is exactly $ab\relu(x)$ when $b>0$. Either way the function depends on the two weights through their product near the point where the unit is switched off, and the unit's parameters fold just as ours do. Networks have many more folds: permuting the hidden units of a layer, flipping the signs of a $\tanh$ unit's weights, rescaling the two sides of a ReLU, and switching units off all produce different weights with the same function \parencite{Sussmann1992,FukumizuAmari2000}. In the language of statistics, neural networks are \term{singular models}\index{singular model}: their parameters are not identifiable, and their Fisher information is degenerate at the parameters that matter \parencite{Watanabe2009,WeiEtAl2023}. The formulas of the regular theory, and the intuitions built on them, such as counting parameters to measure complexity, do not apply.

The theory that does apply was built by Sumio Watanabe from 1999 on \parencite{Watanabe1999,Watanabe2001a,Watanabe2009,Watanabe2018}. It computes the learning coefficient from the geometry of the loss, with the help of a deep theorem of algebraic geometry, Hironaka's resolution of singularities \parencite{Hironaka1964a}, and it proves that this one number, with its multiplicity, replaces $\tfrac d2$ in the free energy and the generalization error. Layered neural networks were among its examples from the first papers \parencite{Watanabe2001a,Watanabe2001b}. Since about 2020 the theory has been applied to modern deep learning \parencite{WeiEtAl2023}: the learning coefficient can be estimated locally at a trained network from samples of a localized posterior \parencite{LauEtAl2025}, and its changes during training have been used to locate stages of development in small transformers \parencite{HooglandEtAl2025}. How much of this carries over to networks trained by stochastic gradient descent, often with more parameters than examples, is a real question, and the last part of this book treats it as one.

\section*{What this book does}
\addcontentsline{toc}{section}{What this book does}

Watanabe's theory has a reputation for being impossible to learn, and the reputation is earned: his books open with algebraic geometry and empirical processes on manifolds, and the connections to deep learning are scattered across papers that assume the theory. This book takes a different road. It starts from $y=abx$ and a handful of relatives, \cref{tab:prologue:models}, on which every concept of the theory can be computed by hand: the free energy, the posterior, the learning coefficient and its multiplicity, the generalization error, the criteria WAIC and WBIC. It adds machinery only when a computation demands it. The one piece of machinery it does not build is resolution of singularities: Hironaka's theorem is stated precisely in the appendix and used as a black box, and every computation of a learning coefficient in the book is done by explicit changes of coordinates that the reader can check.

\begin{table}[t]
\centering
\small
\begin{tabular}{@{}llccc@{}}
\toprule
name & model & $\lambda$ & $m$ & $\tfrac d2$\\
\midrule
the regular plane & $y=ax_1+bx_2$ & $1$ & $1$ & $1$\\
the product model & $y=abx$ & $\tfrac12$ & $2$ & $1$\\
the cusp model & $y=ab^2x$ & $\tfrac14$ & $1$ & $1$\\
the quartic model & $y=ax_1+b^2x_2$ & $\tfrac34$ & $1$ & $1$\\
the chain & $y=a_1a_2\cdots a_Lx$ & $\tfrac12$ & $L$ & $\tfrac L2$\\
the tanh neuron & $y=a\tanh(bx)$ & $\tfrac12$ & $2$ & $1$\\
the ReLU neuron & $y=a\relu(bx)$ & $\tfrac12$ & $2$ & $1$\\
the mixture & $(1-a)\Normal(0,1)+a\Normal(b,1)$ & $\tfrac12$ & $2$ & $1$\\
\bottomrule
\end{tabular}
\caption{Some of the running models, at the truth $w_0=0$ (for the mixture, at the truth $\Normal(0,1)$), with inputs $x\sim\Normal(0,1)$, unit noise, and the uniform prior on $[-1,1]^d$ (for the mixture, $a\in[0,1]$). The learning coefficient $\lambda$ and its multiplicity $m$ are proved in \crefrange{ch:singular}{ch:newton}; $\tfrac d2$ is what the regular theory would charge.}\label{tab:prologue:models}
\end{table}

The book has four parts. \emph{Part I} sets the classical theory against the running models (\cref{ch:regular}), shows how they violate its assumptions (\cref{ch:singular}), and computes the free energy of every model whose loss depends on one coefficient, which turns out to contain most of the theory in miniature (\cref{ch:coefficient}). \emph{Part II} is geometry. The learning coefficient is the rate at which the volume of nearly true parameters shrinks (\cref{ch:volume}); it is computed by blowing up the parameter space until the loss becomes a monomial (\cref{ch:blowups}), or, for many losses, read off a Newton polygon in one step (\cref{ch:newton}); and its values are known for a few families of models in use (\cref{ch:catalogue}). \emph{Part III} is statistics: Watanabe's theorems on the free energy (\cref{ch:freeenergy}) and the generalization error (\cref{ch:generalization}), the criteria WAIC (\cref{ch:waic}) and WBIC (\cref{ch:temperature}), and the way the posterior chooses between singularities (\cref{ch:phases}). \emph{Part IV} turns to trained networks: the local learning coefficient and how to estimate it (\cref{ch:llc}), its changes during training (\cref{ch:stages}), and the gap between Bayesian posteriors and stochastic gradient descent (\cref{ch:sgd}).

Every chapter has the same anatomy. \emph{The question} says what the chapter is about, and a slate box says what the regular theory answers. \emph{By hand} computes the singular answer on one or more of the running models, in numbered steps the reader can check. \emph{The theorem} proves the central result completely; when the general proof needs machinery the book does not build, the chapter proves the theorem under a stated simplifying hypothesis, states the general theorem separately with its source, and says exactly what the general proof adds. \emph{Measuring it} runs a short script, available with the book, and compares what it measures with both theories: the singular prediction is drawn in bronze, the regular one in slate. \emph{The ledger} records the result on a card, and the epilogue collects the fifteen cards. \emph{Where the theory stops} names what was assumed, where the result fails, and what is open. Exercises and notes close the chapter; the notes translate the notation of the main sources into the book's.

\paragraph{How to read.} The chapters are written to be read in order, but each defines what it uses and points to where its tools come from. A short course for a reader who wants the main line and the connection to deep learning is \crefrange{ch:regular}{ch:volume}, then \cref{ch:freeenergy,ch:generalization,ch:llc,ch:sgd}. A reader who wants only the geometry can read Part II after \cref{ch:coefficient}. The prerequisites are multivariable calculus, linear algebra, and probability at the level of a first graduate course: laws of large numbers, the central limit theorem, and convergence in law. No algebraic geometry is assumed.

\paragraph{What this book does not do.} It does not prove Hironaka's theorem, and it does not reproduce the empirical-process theory that Watanabe needs for his theorems in full generality; it states both precisely and proves the theorems completely where they can be proved by elementary means. It treats models with finitely many parameters as the number of examples grows, which is the setting of the theory; it says at each step what that leaves out of deep learning. And it does not claim that singular learning theory explains why neural networks generalize. It explains why parameter counts are the wrong measure of a singular model's complexity in Bayesian learning, and what the right measure is; how far that goes toward explaining deep learning is the subject of \cref{ch:sgd}.

\section*{Notes}
\addcontentsline{toc}{section}{Notes}

The Bayesian information criterion is due to \textcite{Schwarz1978}, who derived it for exponential families; \textcite{Haughton1988} extended it to curved families and proved its consistency. Akaike's criterion, its frequentist counterpart, is due to \textcite{Akaike1974}. The free energy formula for singular models, with the learning coefficient in place of $\tfrac d2$, is Watanabe's: it appears first in \textcite{Watanabe1999} and in full in \textcite{Watanabe2001a}, and the books \parencite{Watanabe2009,Watanabe2018} develop it with its consequences. Resolution of singularities in characteristic zero is the theorem of \textcite{Hironaka1964a}. The symmetries of networks with $\tanh$ units were classified by \textcite{Sussmann1992}; the effect of the resulting singularities on learning dynamics was studied by \textcite{FukumizuAmari2000}. The argument that the singularity of neural networks is a feature for Bayesian learning, not a defect, is made in \textcite{WeiEtAl2023}, and the local learning coefficient as a measure of a trained network's complexity is introduced in \textcite{LauEtAl2025}.

\printchapterbib
\end{refsection}
\sltfloatrestore
